% Previous experiments preserved before nine-page revision (2026-09-26).
% \section{Experiments}
% 
% \subsection{Implement details}
% % text tokenizer
% We train the our model based on Cosmos3 NANO (16B) model~\citep{agarwal2026cosmos}, which consists of a Qwen3-VL(8B) text backbone~\citep{bai2025qwen3}, 36 MoT layers and a Wan2.2 causal VAE~\citep{wan2025wan}. The motion encoder and decoder are trained from scratch while all remaining weights are initialized from the base model. We train the network for 20,000 iterations using the AdamW optimizer with a learning rate of $5 \times 10^{-5}$, governed by a cosine decay schedule with 200 warmup steps. We randomly discard the text embeddings with a probability of 10$\%$ to achieve classifier-free guidance(CFG). Training is conducted in bfloat16(bf16) precision with a bath size of 16 across 16 NVIDIA H100 GPUs. The total training time is approximately 44h. 
% 
% \subsection{Dataset and Evaluation Metrics}
% Our dataset consists of two two parts: one is derived from 3D human-interaction dataset~\citep{li2023omomo, jiang2023chairs, hassan2021samp} and the other is collected from the internet. For the interaction datasets, we render the 3D human motions to obtain human animations. We then use the MiniMax-H3~\citep{minimaxh3} to generate the RGB videos from rendered animations. For the videos downloaded from the internet, we use cameraHMR~\citep{patel2025camerahmr} to obtain the human motions. All RGB videos are annotated by the gemini 2.5~\citep{comanici2025gemini}. Besides, we curate a test dataset from~\citet{li2023omomo, hassan2021samp}, COUCH~\citep{zhang2022couch} and internet to compare with other methods.
% 
% We evaluate the methods from multiply perspectives: 
% 
% \textbf{Motion Quality}: For Motion-prompt alignement, we follow the evaluation protocol defined by MoMask~\citep{guo2024momask} to evaluate the motion-prompt alignment using Frechet Inception Distance(FID), R-Precision(RP)  and MultiModal Distance (MM-Dist). For motion quality, we use Jerk and Foot Sliding Rate(FSR) followed from WHAM~\citep{shin2024wham} and OmniControl to evaluate the jitter~\citep{xie2024omnicontrol} respectively. Since we produce the videos and motions synchronously, we recover the motions from the videos using cameraHMR~\citep{patel2025camerahmr} and calculate the PoseError between extracted motions and generated motions to evaluate the consistency. 
% 
% \textbf{Scene Awareness}: Follow from OMOMO~\citep{li2023omomo}, we adopt contact-F1 as metrics. Besides, we render the generated motions in the corresponding 3D scenes and use QWen3-VL~\citep{bai2025qwen3} to determine whether the interaction is successfully completed. We define the ratio of scuccessful interaction cases over all test cases as the Interaction Success Rate(ISR).
% 
% \textbf{Speed}: We record the mean inference time of each method for Speed comparisons.
% 
% 
% \subsection{Comparison Results}
% We compare our method with three types of method: T2M methods, HOI methods and the I2V+HMR method. For T2M methods, we compare with Kimodo~\citep{Kimodo2026}, HY-Motion~\citep{hymotion2025} and Go-to-Zero~\citep{fan2025gotozero}. For HOI methods, we compare with LIGHT~\citep{wang2026light} and HOIFHLI~\citep{wu2025hoifhli}. For the two-stage inference method, we leverage MiniMax-H3~\citep{minimaxh3} to generate videos and CameraHMR~\citep{patel2025camerahmr} for subsequent motion recovery. CoMoVi~\citep{zhao2026comovi} is the most related work to ours. However, it additionally relies on a 2D motion representation for generation, whereas our method trains with less data requirements. we also incorporate CoMoVi~\citep{zhao2026comovi} as our baseline. For fair comparisons, the methods that conditioned on an initial human pose(Kimodo, HOIFHLI, CoMoVi, Ours) use the motion extracted by cameraHMR, which may introduce some space errors(Figure~\ref{fig:Comparison}(a)). In addition, we fine-tune CoMoVi on our training dataset for comparison. Since HOI methods generate human motions and object motions simultaneously, we will visualize the object trajectories generated by HOIFHLI and LIGHT. The complete quantitative results and qualitative results are shown in Table~\ref{tab:comparison} and Figure~\ref{fig:Comparison}.
% 
% \begin{table}[!htbp]
% \caption{Quantitative results of motion generation methods.}
% \label{tab:comparison}
% \begin{center}
% \resizebox{\textwidth}{!}{%
% \begin{tabular}{lccccccccc}
% \toprule
% \multirow{2}{*}{\textbf{Method}}
% & \multicolumn{3}{c}{\textbf{Motion-Prompt Alignment}}
% & \multicolumn{2}{c}{\textbf{Scene Awareness}}
% & \multicolumn{2}{c}{\textbf{Jitter}}
% & \multirow{2}{*}{\textbf{PoseError} $\downarrow$}
% & \multirow{2}{*}{\textbf{Time} $\downarrow$} \\
% \cmidrule(lr){2-4}
% \cmidrule(lr){5-6}
% \cmidrule(lr){7-8}
% & FID $\downarrow$
% & RP@1 $\uparrow$
% & MM-Dist $\downarrow$
% & Contact-F1 $\uparrow$
% & ISR $\uparrow$
% & Jerk $\downarrow$
% & FSR $\downarrow$
% & & \\
% 
% \midrule
% \textbf{T2M methods} & & & & & & & & & \\
% 
% \hspace{1em}HY-Motion
% & 0.46
% & 0.10
% & 1.21
% & 0.09
% & 0.51
% & 3.43
% & \textbf{1.25}
% & --
% & 2.83s \\
% 
% \hspace{1em}Kimodo
% & 0.66
% & 0.05
% & 1.30
% & 0.03
% & 0.23
% & 35.14
% & 6.07
% & --
% & \textbf{1.93s} \\
% 
% \hspace{1em}Go-to-Zero
% & 0.51
% & 0.04
% & 1.35
% & 0.07
% & 0.36
% & 17.50
% & 3.42
% & --
% & 2.57s \\
% 
% \addlinespace[0.5em]
% \multicolumn{10}{l}{\textbf{HOI methods}} \\
% 
% \hspace{1em}HOIFHLI
% & 0.82
% & 0.02
% & 1.37
% & 0.09
% & 0.12
% & 10.62
% & 6.28
% & --
% & 13.94s \\
% 
% \hspace{1em}LIGHT
% & 0.62
% & 0.03
% & 1.35
% & 0.17
% & 0.46
% & \textbf{1.37}
% & 11.60
% & --
% & 25.88s \\
% 
% \addlinespace[0.5em]
% \textbf{I2V+HMR} & & & & & & & & & \\
% 
% \begin{tabular}[c]{@{}l@{}}
% \hspace{1em}MiniMax-H3 +\\
% \hspace{1em}CameraHMR
% \end{tabular}
% & 0.69
% & 0.10
% & 1.32
% & \textbf{0.24}
% & \textbf{0.91}
% & 37.22
% & 6.41
% & --
% & 523.67s \\
% 
% \addlinespace[0.5em]
% \multicolumn{10}{l}{\textbf{Baseline}} \\
% 
% \hspace{1em}CoMoVi
% & 0.71
% & 0.12
% & 1.25
% & 0.09
% & 0.69
% & 16.10
% & 24.45
% & 81.28
% & 270.45s \\
% 
% \midrule
% 
% \hspace{1em}\textbf{Ours}
% & \textbf{0.27}
% & \textbf{0.24}
% & \textbf{1.13}
% & 0.22
% & \textbf{0.91}
% & 8.35
% & 2.80
% & \textbf{67.50}
% & 156.59s \\
% 
% \bottomrule
% \end{tabular}%
% }
% \end{center}
% \end{table}
% 
% \begin{figure}[!h]
%     \centering
%     \includegraphics[width=\linewidth]{Fig/pdf/Comparison.pdf}
%     \caption{Qualitative comparison of motion generation methods.}
%     \label{fig:Comparison}
% \end{figure}
% 
% \textbf{T2M methods.} Our method significantly outperforms T2M methods in both Motion-Prompt alignment and scene-awareness. Due to the lack of scene awareness, T2M methods fail to sit at the intended locations. T2M methods generally achieve lower jitter and FSR by encoding the velocity and acceleartion into motion representations. Since the motion stream contains fewer tokens than the video stream, T2M methods nearly achieve real-time inference.
% 
% \textbf{HOI methods.} Both HOIFHLI and LIGHT rely on 3D object information as input. LIGHT additionally requires human shape, while HOIFHLI requires the initial human pose. Our method generally outperforms the HOI methods in motion-prompt alignment and jitter. Since LIGHT and HOIFHLI are mainly optimized for dynamic interaction, Our method achieve better scene awareness as shown in Figure~\ref{fig:Comparison}(b)(c). Since the HOI methods don't need to denoise video latents, they are faster than our methods.
% 
% \textbf{I2V+HMR.} Our method outperforms the MiniMax-H3+CameraHMR pipeline across all evaluated metrics except Contact-F1. Because jointly generating introduces a trade-off in video quality, which further affect interaction accuracy in a little cases. In particular, the two-stage pipeline requires about 3$\times$ longer inference time.
% 
% \textbf{Baseline.} Our method also achieve better performance than the baseline across all metrics. CoMoVi predicts all joint rotations without root translations, making it effective for motions with limited displacement(Figure\ref{fig:Comparison}(a)). When the human needs to move, the motion accuracy drops significantly(Figure~\ref{fig:Comparison}(b)(c)).
% 
% 
% \subsection{Abalation Study}
% To verify that our strategies effectively reduce jitter, we conduct two ablation studies to evaluate the effectiveness. We use Jerk and Foot Sliding Rate as metrics to quantify the jitter. For a direct comparisons, we don't apply zero-phase Savitzky-Golay filter to any experimental settings.
% 
% \textbf{Motion Representation.} To validate the effectiveness of our motion representation, we use the absolute translation of the root node for comparison. We train the model with two motion representations while maintaining other settings fixed. The results are shown in Table\ref{tab:ablation_norm}. The experiment shows that relative translation representation outperforms the absolute translation representation on both metrics. Directly modeling absolute translations introduces noisy in the global coordinate space. The errors are magnified when the predictions are rescaled back to the original coordinate range, resulting in more severe temporal jitter.
% % \begin{table}[h]
% % \caption{Ablation study on motion representation.}
% % \label{tab:ablation_norm}
% % \begin{center}
% % \begin{tabular}{lccc}
% % \toprule
% %  & Jerk$\downarrow$  & Foot sliding rate $\downarrow$ \\
% % \midrule
% % % w/o relative translation & 4.355 $\pm$ 4.425 & 0.766 & \\
% % % w/ relative translation & \textbf{1.251 $\pm$ 0.973} & \textbf{0.457} & \\
% % w/o relative translation & 4.355 & 0.766 & \\
% % w/ relative translation & \textbf{1.251} & \textbf{0.457} & \\
% % \bottomrule
% % \end{tabular}
% % \end{center}
% % \end{table}
% \begin{table}[h]
% \caption{Root translation ablation without post-processing. Translations are either absolute or relative to the initial root position; all other settings are fixed. Lower values are better.}
% \label{tab:ablation_norm}
% \begin{center}
% \begin{tabular}{lcc}
% \toprule
%  & Jerk $\downarrow$  & Foot sliding rate $\downarrow$ \\
% \midrule
% % w/o relative translation & 4.355 $\pm$ 4.425 & 0.766 & \\
% % w/ relative translation & \textbf{1.251 $\pm$ 0.973} & \textbf{0.457} & \\
% Absolute root translation & 4.355 & 0.766 \\
% Relative root translation & 1.251 & 0.457 \\
% \bottomrule
% \end{tabular}
% \end{center}
% \end{table}
% 
% 
% % \textbf{Coupled Dual-timestep Schedule.}
% \noindent Shift-Decoupled Noise Schedule.
% % TODO: 关于去噪理论的解释，假定每个token包含的信息密度一样，那么数量越多的token越需要进行去噪在高频，从而让去噪的平均值和token数少的保持一致，这点应该要重新再写下
% 
% % Before correcting the denoising figure reference (2026-09-26).
% % To validate the effectiveness of the coupled noise schedule and train-inference consistency, we conduct nine ablation experiments as shown in Table~\ref{tab:ablation_schedule}. We train three shift configurations while keeping other parameters fixed. For each training shift configurations, the lowest jitter is observed when the training and inference shifts are kept consistent. Since the model is trained on noise pairs along a single coupled trajectory (Figure\ref{fig:Method} (c)), Using noise pairs deviated from the trajectory will lead to prediction errors and more severe jitter. Comparing the different shift configurations, we find that \{ $s_v$ = 10, $s_m$ = 1 \} achieve the best performance. Since the motion latents is more compact and less redundancy, the uniform denoising schedule is more sufficient than keeping the latent longer in the high noise region.
% % To validate the effectiveness of the coupled noise schedule and train-inference consistency, we conduct nine ablation experiments as shown in Table~\ref{tab:ablation_schedule}. We train three shift configurations while keeping other parameters fixed. For each training shift configurations, the lowest jitter is observed when the training and inference shifts are kept consistent. Since the model is trained on noise pairs along a single coupled trajectory (Figure~\ref{fig:denoising}(b)), Using noise pairs deviated from the trajectory will lead to prediction errors and more severe jitter. Comparing the different shift configurations, we find that \{ $s_v$ = 10, $s_m$ = 1 \} achieve the best performance. Since the motion latents is more compact and less redundancy, the uniform denoising schedule is more sufficient than keeping the latent longer in the high noise region.
% We evaluate three training shift configurations, each with three inference configurations, while keeping other parameters fixed (Table~\ref{tab:ablation_schedule}). For every training configuration, matching the inference shifts to the training shifts gives the lowest jerk and foot sliding rate. Among all tested configurations, matched training and inference shifts of $s_v=10$ and $s_m=1$ achieve the lowest values on both metrics. These results support using modality-specific shifts and maintaining the same noise relationship during training and inference.
% 
% 
% % \begin{table}[!htbp]
% % \caption{Ablation study on Coupled Dual-timestep Schedule.}
% % \label{tab:ablation_schedule}
% % \begin{center}
% % \begin{tabular}{ccccc}
% % \toprule
% % Training & Inference & Jerk$\downarrow$ & Foot Sliding ratio$\downarrow$ \\
% % \midrule
% % % \multirow{3}{*}{$s_v=10,s_m=1$}
% % %  & $s_v=10, s_m=1$ & \textbf{1.251 $\pm$ 0.97} & \textbf{0.457} \\
% % %  & $s_v=10, s_m=10$ & 8.003 $\pm$ 2.48 & 0.986 \\
% % %  & $s_v=10, s_m=15$ & 9.703 $\pm$ 3.47 & 1.000\\
% % % \midrule
% % % \multirow{3}{*}{$s_v=10,s_m=10$}
% % %  & $s_v=10, s_m=1$ & 7.021 $\pm$ 1.04 & 0.985\\
% % %  & $s_v=10, s_m=10$ & \textbf{1.635} $\pm$ \textbf{1.05} & \textbf{0.635} \\
% % %  & $s_v=10, s_m=15$ & 2.557 $\pm$ 0.79 & 0.830 \\
% % % \midrule
% % % \multirow{3}{*}{$s_v=10,s_m=15$}
% % %  & $s_v=10, s_m=1$ & 8.413 $\pm$ 1.01 & 0.981\\
% % %  & $s_v=10, s_m=10$ & 2.127 $\pm$ 1.55 & 0.700\\
% % %  & $s_v=10, s_m=15$ & \textbf{1.809} $\pm$ \textbf{1.28} & \textbf{0.649} \\
% % \multirow{3}{*}{$s_v=10,s_m=1$}
% %  & $s_v=10, s_m=1$ & \textbf{1.251} & \textbf{0.457} \\
% %  & $s_v=10, s_m=10$ & 8.003 & 0.986 \\
% %  & $s_v=10, s_m=15$ & 9.703 & 1.000\\
% % \midrule
% % \multirow{3}{*}{$s_v=10,s_m=10$}
% %  & $s_v=10, s_m=1$ & 7.021 & 0.985\\
% %  & $s_v=10, s_m=10$ & \textbf{1.635} & \textbf{0.635} \\
% %  & $s_v=10, s_m=15$ & 2.557 & 0.830 \\
% % \midrule
% % \multirow{3}{*}{$s_v=10,s_m=15$}
% %  & $s_v=10, s_m=1$ & 8.413 & 0.981\\
% %  & $s_v=10, s_m=10$ & 2.127 & 0.700\\
% %  & $s_v=10, s_m=15$ & \textbf{1.809} & \textbf{0.649} \\
% % \bottomrule
% % \end{tabular}
% % \end{center}
% % \end{table}
% \begin{table}[!htbp]
% \caption{Ablation of the shift-decoupled noise schedule without post-processing. Each group fixes the training shifts and varies the inference shifts. Matching shifts yields the lowest jerk and foot sliding rate within each group.}
% \label{tab:ablation_schedule}
% \begin{center}
% \begin{tabular}{cccc}
% \toprule
% Training shifts & Inference shifts & Jerk $\downarrow$ & Foot sliding rate $\downarrow$ \\
% \midrule
% % \multirow{3}{*}{$s_v=10,s_m=1$}
% %  & $s_v=10, s_m=1$ & 1.251 $\pm$ 0.97 & 0.457 \\
% %  & $s_v=10, s_m=10$ & 8.003 $\pm$ 2.48 & 0.986 \\
% %  & $s_v=10, s_m=15$ & 9.703 $\pm$ 3.47 & 1.000\\
% % \midrule
% % \multirow{3}{*}{$s_v=10,s_m=10$}
% %  & $s_v=10, s_m=1$ & 7.021 $\pm$ 1.04 & 0.985\\
% %  & $s_v=10, s_m=10$ & 1.635 $\pm$ 1.05 & 0.635 \\
% %  & $s_v=10, s_m=15$ & 2.557 $\pm$ 0.79 & 0.830 \\
% % \midrule
% % \multirow{3}{*}{$s_v=10,s_m=15$}
% %  & $s_v=10, s_m=1$ & 8.413 $\pm$ 1.01 & 0.981\\
% %  & $s_v=10, s_m=10$ & 2.127 $\pm$ 1.55 & 0.700\\
% %  & $s_v=10, s_m=15$ & 1.809 $\pm$ 1.28 & 0.649 \\
% \multirow{3}{*}{$s_v=10,s_m=1$}
%  & $s_v=10, s_m=1$ & 1.251 & 0.457 \\
%  & $s_v=10, s_m=10$ & 8.003 & 0.986 \\
%  & $s_v=10, s_m=15$ & 9.703 & 1.000\\
% \midrule
% \multirow{3}{*}{$s_v=10,s_m=10$}
%  & $s_v=10, s_m=1$ & 7.021 & 0.985\\
%  & $s_v=10, s_m=10$ & 1.635 & 0.635 \\
%  & $s_v=10, s_m=15$ & 2.557 & 0.830 \\
% \midrule
% \multirow{3}{*}{$s_v=10,s_m=15$}
%  & $s_v=10, s_m=1$ & 8.413 & 0.981\\
%  & $s_v=10, s_m=10$ & 2.127 & 0.700\\
%  & $s_v=10, s_m=15$ & 1.809 & 0.649 \\
% \bottomrule
% \end{tabular}
% \end{center}
% \end{table}
% 
% 
% 
% \label{sec:experiments}
% 
% 

% Previous nine-page experiments preserved before structure and table revision.
% \clearpage
% \section{Experiments}
% \label{sec:experiments}
% \subsection{Experimental Setup}
% \label{sec:experimental-setup}
% 
% \noindent Training. We initialize World2Motion from Cosmos 3 NANO (16B)~\citep{agarwal2026cosmos}, which includes a Qwen3-VL (8B) text backbone~\citep{bai2025qwen3}, 36 Mixture-of-Transformers layers, and a Wan2.2 causal VAE~\citep{wan2025wan}. The motion encoder and decoder are trained from scratch, and the understanding branch remains frozen. We train for 20,000 iterations with AdamW, a learning rate of $5\times10^{-5}$, a cosine decay schedule, and 200 warmup steps. Text conditioning is dropped with probability 0.1 for classifier-free guidance. Training uses bfloat16 precision, a batch size of 16, and 16 NVIDIA H100 GPUs, and takes approximately 44 hours.
% 
% \noindent Training and test data. The hybrid training data combines synthetic video--motion pairs derived from OMOMO~\citep{li2023omomo}, CHAIRs~\citep{jiang2023chairs}, and SAMP~\citep{hassan2021samp} with internet videos paired with CameraHMR~\citep{patel2025camerahmr} motion estimates. Section~\ref{sec:training-data} describes their construction. The synthetic component supplies known 3D motion, while the real component expands the range of environments and appearances. Text descriptions are generated using Gemini 2.5~\citep{comanici2025gemini}. We evaluate on a test set assembled from OMOMO, SAMP, COUCH~\citep{zhang2022couch}, and internet videos. This multi-source evaluation tests whether generated motions follow the requested action and its spatial relationship to the depicted environment.
% 
% \noindent Baselines and conditioning. We compare with text-to-motion methods HY-Motion~\citep{hymotion2025}, Kimodo~\citep{Kimodo2026}, and Go-to-Zero~\citep{fan2025gotozero}; human--object interaction methods LIGHT~\citep{wang2026light} and HOIFHLI~\citep{wu2025hoifhli}; and a two-stage MiniMax-H3~\citep{minimaxh3} + CameraHMR pipeline. We also fine-tune CoMoVi~\citep{zhao2026comovi} on our training data as a joint video--motion baseline. Methods that accept an initial human pose---Kimodo, HOIFHLI, CoMoVi, and World2Motion---use the pose estimated by CameraHMR. This shared pose estimate can contain spatial errors that affect the subsequent motion. LIGHT and HOIFHLI additionally require explicit 3D object information; World2Motion obtains scene context from the input image. For the interaction baselines, we visualize their predicted object trajectories together with the human motion.
% 
% \noindent Motion distribution and text alignment. We follow the MoMask~\citep{guo2024momask} evaluation protocol. Fr\'echet Inception Distance (FID) measures the difference between generated and reference motion feature distributions. R-Precision at rank 1 (R@1) and multimodal distance (MM-Dist) evaluate motion--text alignment. Lower FID and MM-Dist and higher R@1 are better. These scores assess different properties: a motion can match an action description while failing to reach the intended object.
% 
% \noindent Scene interaction and motion stability. Following OMOMO, we report Contact-F1. We also render generated motions in the corresponding 3D scenes and use Qwen3-VL to assess whether the requested interaction is completed. Interaction Success Rate (ISR) is the fraction of test cases judged successful. Jerk measures temporal motion irregularity, and foot sliding rate (FSR) measures sliding artifacts, following the evaluation procedures of WHAM~\citep{shin2024wham} and OmniControl~\citep{xie2024omnicontrol}. Lower values are better for both stability metrics.
% 
% \noindent Video--motion consistency and efficiency. For joint generation methods, we recover motion from the generated video using CameraHMR and report PoseError against the directly generated motion. This evaluates agreement between the two outputs. We also report mean inference time in seconds. World2Motion results in the main comparison use the post-processing described in Section~\ref{sec:method}; all ablation results are obtained without the Savitzky--Golay filter so that representation and scheduling choices can be assessed directly.
% 
% \noindent Interpreting the metrics. PoseError depends on video-based motion recovery and measures cross-modal agreement rather than error against a ground-truth 3D trajectory. We therefore read it alongside text alignment and interaction scores. Similarly, low jerk does not establish that a requested interaction is completed: smooth motion can miss an object, while a successful interaction may still contain local artifacts. Reporting these measures together separates scene interaction from temporal stability and provides a more complete assessment of scene-aware motion generation.
% 
% \clearpage
% \subsection{Quantitative Comparisons}
% 
% \begin{table}[!htbp]
% \centering
% \small
% \setlength{\tabcolsep}{4pt}
% \caption{Comparison on the multi-source test set. (a) Motion distribution, text alignment, and scene interaction. (b) Motion stability, video--motion consistency, and mean inference time. Arrows indicate better values; underlining marks the best reported value in each column. Dashes indicate unreported scores.}
% \label{tab:comparison}
% \begin{tabular*}{\linewidth}{@{\extracolsep{\fill}}lccccc@{}}
% \toprule
% \multicolumn{6}{l}{(a) Motion generation and scene interaction}\\
% Method & FID $\downarrow$ & R@1 $\uparrow$ & MM-Dist $\downarrow$ & Contact-F1 $\uparrow$ & ISR $\uparrow$\\
% \midrule
% \multicolumn{6}{l}{Text-to-motion}\\
% HY-Motion & 0.46 & 0.10 & 1.21 & 0.09 & 0.51\\
% Kimodo & 0.66 & 0.05 & 1.30 & 0.03 & 0.23\\
% Go-to-Zero & 0.51 & 0.04 & 1.35 & 0.07 & 0.36\\
% \midrule
% \multicolumn{6}{l}{Human--object interaction}\\
% HOIFHLI & 0.82 & 0.02 & 1.37 & 0.09 & 0.12\\
% LIGHT & 0.62 & 0.03 & 1.35 & 0.17 & 0.46\\
% \midrule
% \multicolumn{6}{l}{Two-stage video generation and motion recovery}\\
% MiniMax-H3 + CameraHMR & 0.69 & 0.10 & 1.32 & \underline{0.24} & \underline{0.91}\\
% \midrule
% \multicolumn{6}{l}{Joint video--motion generation}\\
% CoMoVi & 0.71 & 0.12 & 1.25 & 0.09 & 0.69\\
% World2Motion & \underline{0.27} & \underline{0.24} & \underline{1.13} & 0.22 & \underline{0.91}\\
% \bottomrule
% \end{tabular*}
% 
% \medskip
% \begin{tabular*}{\linewidth}{@{\extracolsep{\fill}}lcccc@{}}
% \toprule
% \multicolumn{5}{l}{(b) Motion stability, output consistency, and efficiency}\\
% Method & Jerk $\downarrow$ & FSR $\downarrow$ & PoseError $\downarrow$ & Time (s) $\downarrow$\\
% \midrule
% \multicolumn{5}{l}{Text-to-motion}\\
% HY-Motion & 3.43 & \underline{1.25} & -- & 2.83\\
% Kimodo & 35.14 & 6.07 & -- & \underline{1.93}\\
% Go-to-Zero & 17.50 & 3.42 & -- & 2.57\\
% \midrule
% \multicolumn{5}{l}{Human--object interaction}\\
% HOIFHLI & 10.62 & 6.28 & -- & 13.94\\
% LIGHT & \underline{1.37} & 11.60 & -- & 25.88\\
% \midrule
% \multicolumn{5}{l}{Two-stage video generation and motion recovery}\\
% MiniMax-H3 + CameraHMR & 37.22 & 6.41 & -- & 523.67\\
% \midrule
% \multicolumn{5}{l}{Joint video--motion generation}\\
% CoMoVi & 16.10 & 24.45 & 81.28 & 270.45\\
% World2Motion & 8.35 & 2.80 & \underline{67.50} & 156.59\\
% \bottomrule
% \end{tabular*}
% \end{table}
% 
% \noindent Motion generation and scene interaction. World2Motion achieves the lowest FID and MM-Dist and the highest R@1 among the compared methods (Table~\ref{tab:comparison}(a)). Relative to HY-Motion, R@1 increases from 0.10 to 0.24 and ISR from 0.51 to 0.91. Its ISR exceeds those of the evaluated text-to-motion and human--object interaction methods. The two-stage baseline reaches the same ISR and a higher Contact-F1 (0.24 versus 0.22). Thus, the results support improved action alignment and competitive scene interaction, while leaving room to improve contact accuracy.
% 
% \noindent Stability and output consistency. World2Motion has lower jerk and FSR than Kimodo, Go-to-Zero, HOIFHLI, and the two video-based baselines. HY-Motion achieves lower values on both stability metrics, and LIGHT achieves lower jerk. World2Motion therefore does not lead every measure of motion quality. Compared with CoMoVi, it reduces PoseError from 81.28 to 67.50, a 17.0\% reduction, indicating closer agreement between the generated video and 3D motion.
% 
% \noindent Inference efficiency. World2Motion takes 156.59 seconds on average, compared with 523.67 seconds for MiniMax-H3 + CameraHMR and 270.45 seconds for CoMoVi: speedups of approximately $3.34\times$ and $1.73\times$, respectively. The text-to-motion and human--object interaction methods remain faster. Our efficiency gain concerns video-based motion generation; jointly producing video and motion still costs more than generating motion alone.
% 
% \clearpage
% \subsection{Qualitative Comparisons}
% 
% \begin{figure}[!htbp]
% \centering
% \includegraphics[width=0.96\linewidth]{Fig/pdf/Comparison.pdf}
% \caption{Qualitative comparison of scene interactions: (a) lifting a plastic box, (b) walking to a chair and sitting, and (c) walking to a sofa and sitting. Each method shows successive motion states; video-generating methods also show corresponding video frames. Human--object interaction baselines include their predicted object trajectories.}
% \label{fig:Comparison}
% \end{figure}
% 
% Figure~\ref{fig:Comparison} compares actions that depend on both body motion and placement relative to an object. In the sitting examples, producing a plausible sitting pose is insufficient: the character must first approach the intended seat and then align the body with it. The box example also involves coordinating the body with object motion. Viewing the approach and final interaction together helps distinguish action matching from scene-consistent execution. For video-based methods, the paired rows additionally allow the reader to inspect whether the generated video and recovered or directly generated 3D motion describe the same interaction.
% 
% \clearpage
% \subsection{Ablation Studies}
% 
% We evaluate root translation and noise scheduling using jerk and FSR. All ablations disable post-processing. Comparisons within each study keep other settings fixed. The first study compares motion representations; the second changes training and inference shifts independently. This separates the choice of noise schedule from the effect of using it consistently.
% 
% \noindent Root translation. Representing each root position relative to the initial position reduces jerk from 4.355 to 1.251 (71.3\%) and FSR from 0.766 to 0.457 (40.3\%), compared with absolute translation (Table~\ref{tab:ablation_norm}). Both representations retain root position and 6D joint rotations; the change is the reference position used for translation. The comparison therefore isolates this coordinate choice while keeping the rotation representation and other settings fixed. These results support using relative positions for more stable generation without requiring root velocity as an input.
% 
% \begin{table}[!htbp]
% \centering
% \small
% \setlength{\tabcolsep}{8pt}
% \caption{Root translation ablation without post-processing. Relative positions are measured from the initial root position. Underlining marks the lower value in each column.}
% \label{tab:ablation_norm}
% \begin{tabular}{lcc}
% \toprule
% Root translation & Jerk $\downarrow$ & FSR $\downarrow$\\
% \midrule
% Absolute & 4.355 & 0.766\\
% Relative to initial position & \underline{1.251} & \underline{0.457}\\
% \bottomrule
% \end{tabular}
% \end{table}
% 
% \noindent Shift-decoupled noise schedule. Table~\ref{tab:ablation_schedule} separates two questions: which shifts to use, and whether the training and inference shifts should match. With matched shifts, $s_m=1$ achieves lower jerk and FSR than $s_m=10$ or $15$. Compared with matched $s_m=10$, matched $s_m=1$ reduces jerk by 23.5\% and FSR by 28.0\%. For each training configuration, matching the inference shift also gives the lowest values on both metrics. For example, after training with $s_m=1$, changing only the inference shift to $10$ increases jerk from 1.251 to 8.003. These observations support modality-specific shifts and consistent training and inference noise relationships. They do not require an additional smoothing step, since the filter is disabled throughout this comparison.
% 
% \begin{table}[!htbp]
% \centering
% \small
% \setlength{\tabcolsep}{10pt}
% \caption{Noise-schedule ablation without post-processing. Video shift $s_v=10$ is fixed in training and inference. Groups vary the training motion shift; underlining marks the best result within each group.}
% \label{tab:ablation_schedule}
% \begin{tabular}{cccc}
% \toprule
% Training $s_m$ & Inference $s_m$ & Jerk $\downarrow$ & FSR $\downarrow$\\
% \midrule
% \multirow{3}{*}{1} & 1 & \underline{1.251} & \underline{0.457}\\
%  & 10 & 8.003 & 0.986\\
%  & 15 & 9.703 & 1.000\\
% \midrule
% \multirow{3}{*}{10} & 1 & 7.021 & 0.985\\
%  & 10 & \underline{1.635} & \underline{0.635}\\
%  & 15 & 2.557 & 0.830\\
% \midrule
% \multirow{3}{*}{15} & 1 & 8.413 & 0.981\\
%  & 10 & 2.127 & 0.700\\
%  & 15 & \underline{1.809} & \underline{0.649}\\
% \bottomrule
% \end{tabular}
% \end{table}
% \FloatBarrier
% 

% Allow Figure 3 to share a page with the following text.
\section{Experiments}
\label{sec:experiments}
\subsection{Experimental Setup}
\label{sec:experimental-setup}

% Before metric-protocol clarification: \noindent\textbf{Dataset.} Our training dataset combines synthetic video--motion pairs from OMOMO~\citep{li2023omomo}, CHAIRs~\citep{jiang2023chairs}, and SAMP~\citep{hassan2021samp} with internet videos paired with CameraHMR~\citep{patel2025camerahmr} motion estimates. Section~\ref{sec:training-data} describes their construction. The synthetic component supplies known 3D motion, while the real component expands the range of environments and appearances. Text descriptions are generated using Gemini 2.5~\citep{comanici2025gemini}. We evaluate on a test set assembled from OMOMO, SAMP, COUCH~\citep{zhang2022couch}, and internet videos.
% Before dataset-and-splits revision: \noindent\textbf{Dataset.} Our training dataset combines synthetic video--motion pairs from OMOMO~\citep{li2023omomo}, CHAIRs~\citep{jiang2023chairs}, and SAMP~\citep{hassan2021samp} with internet videos paired with CameraHMR~\citep{patel2025camerahmr} motion estimates. Section~\ref{sec:training-data} describes their construction. The synthetic component supplies known 3D motion, while the real component expands the range of environments and appearances. Text descriptions are generated using Gemini 2.5~\citep{comanici2025gemini}. The main comparison uses 100 interaction cases: 12 from CHAIRs, 70 from COUCH~\citep{zhang2022couch}, 12 from OMOMO, and 6 from SAMP. The ablations use a separate 35-case set, described in Section~\ref{sec:ablation}.
% Before experiment citation revision: \noindent\textbf{Dataset Construction and Splits.} Following Section~\ref{sec:training-data}, HOI-mix combines SAMP, CHAIRs, and OMOMO motions paired with synthesized RGB videos and curated Bilibili videos with CameraHMR motion estimates. The training set contains 20,812 video--motion--text triplets totaling 29.15 hours, each with 121 frames at 24 fps. These triplets represent 10,847 distinct motion clips, with two appearance variants per interaction motion and one video per Bilibili clip. Motions share a common camera coordinate system; training uses sample-level shuffling across the merged dataset without source-specific reweighting. The legacy validation split contains 30 interaction and 20 non-interaction clips. Interaction data are held out by source sequence, and Bilibili data by clip. The final interaction benchmark contains 100 cases: 12 from CHAIRs, 12 from OMOMO, 6 from SAMP, and 70 additional cases from COUCH~\citep{zhang2022couch}. All benchmark source sequences are disjoint from the current training manifest. However, the benchmark retains the 30 interaction cases from the legacy validation split and is therefore not independent of that split. The separate 35-case ablation set is described in Section~\ref{sec:ablation}.
\noindent\textbf{Dataset Construction and Splits.} Following Section~\ref{sec:training-data}, HOI-mix combines SAMP~\citep{hassan2021samp}, CHAIRs~\citep{jiang2023chairs}, and OMOMO~\citep{li2023omomo} motions paired with synthesized RGB videos and curated Bilibili videos with CameraHMR~\citep{patel2025camerahmr} motion estimates. The training set contains 20,812 video--motion--text triplets totaling 29.15 hours, each with 121 frames at 24 fps. These triplets represent 10,847 distinct motion clips, with two appearance variants per interaction motion and one video per Bilibili clip. Motions share a common camera coordinate system; training uses sample-level shuffling across the merged dataset without source-specific reweighting. The legacy validation split contains 30 interaction and 20 non-interaction clips. Interaction data are held out by source sequence, and Bilibili data by clip. The final interaction benchmark contains 100 cases: 12 from CHAIRs, 12 from OMOMO, 6 from SAMP, and 70 additional cases from COUCH~\citep{zhang2022couch}. All benchmark source sequences are disjoint from the current training manifest. However, the benchmark retains the 30 interaction cases from the legacy validation split and is therefore not independent of that split. The separate 35-case ablation set is described in Section~\ref{sec:ablation}.

% Before experiment citation revision: \noindent\textbf{Implementation Details.} We initialize World2Motion from Cosmos 3 NANO (16B)~\citep{agarwal2026cosmos}, which includes a Qwen3-VL (8B) text backbone~\citep{bai2025qwen3}, 36 Mixture-of-Transformers layers, and a Wan2.2 causal VAE~\citep{wan2025wan}. The motion encoder and decoder are trained from scratch, and the understanding branch remains frozen. We train for 20,000 iterations with AdamW, a learning rate of $5\times10^{-5}$, a cosine decay schedule, and 200 warmup steps. Text conditioning is dropped with probability 0.1 for classifier-free guidance. Training uses bfloat16 precision, a batch size of 16, and 16 NVIDIA H100 GPUs, and takes approximately 44 hours.
\noindent\textbf{Implementation Details.} We initialize World2Motion from Cosmos 3 NANO (16B)~\citep{agarwal2026cosmos}, which includes a Qwen3-VL (8B) text backbone~\citep{bai2025qwen3}, 36 Mixture-of-Transformers layers, and a Wan2.2 causal VAE~\citep{wan2025wan}. The motion encoder and decoder are trained from scratch, and the understanding branch remains frozen. We train for 20,000 iterations with AdamW~\citep{loshchilov2019decoupled}, a learning rate of $5\times10^{-5}$, a cosine decay schedule, and 200 warmup steps. Text conditioning is dropped with probability 0.1 for classifier-free guidance~\citep{ho2022classifierfree}. Training uses bfloat16 precision, a batch size of 16, and 16 NVIDIA H100 GPUs, and takes approximately 44 hours.

% Before dataset-and-splits revision: \noindent\textbf{Baselines.} We compare with text-to-motion methods HY-Motion~\citep{hymotion2025}, Kimodo~\citep{Kimodo2026}, and Go-to-Zero~\citep{fan2025gotozero}; human--object interaction methods LIGHT~\citep{wang2026light} and HOIFHLI~\citep{wu2025hoifhli}; and a two-stage MiniMax-H3~\citep{minimaxh3} + CameraHMR pipeline. We also fine-tune CoMoVi~\citep{zhao2026comovi} on our training data as a joint video--motion baseline. Methods that accept an initial human pose---Kimodo, HOIFHLI, CoMoVi, and World2Motion---use the pose estimated by CameraHMR. This shared pose estimate can contain spatial errors that affect the subsequent motion. LIGHT and HOIFHLI additionally require explicit 3D object information; World2Motion obtains scene context from the input image.
% Before experiment citation revision: \noindent\textbf{Baselines.} We compare with text-to-motion methods HY-Motion~\citep{hymotion2025}, Kimodo~\citep{Kimodo2026}, and Go-to-Zero~\citep{fan2025gotozero}; interaction methods LIGHT~\citep{wang2026light} and HOIFHLI~\citep{wu2025hoifhli}; MiniMax-H3~\citep{minimaxh3} + CameraHMR; and CoMoVi~\citep{zhao2026comovi}, fine-tuned on our data. Kimodo, HOIFHLI, CoMoVi, and World2Motion use CameraHMR initial poses. LIGHT and HOIFHLI additionally require 3D object information; World2Motion uses the scene image.
\noindent\textbf{Baselines.} We compare with text-to-motion methods HY-Motion~\citep{hymotion2025}, Kimodo~\citep{Kimodo2026}, and Go-to-Zero~\citep{fan2025gotozero}; interaction methods LIGHT~\citep{wang2026light} and HOIFHLI~\citep{wu2025hoifhli}; MiniMax-H3~\citep{minimaxh3} + CameraHMR; and CoMoVi~\citep{zhao2026comovi}, fine-tuned on our data. We initialize Kimodo, HOIFHLI, CoMoVi, and World2Motion with poses estimated by CameraHMR~\citep{patel2025camerahmr}. LIGHT and HOIFHLI additionally require 3D object information, while World2Motion uses the scene image.

\subsection{Evaluation}
\label{sec:evaluation}

% Before metric-protocol clarification: \noindent\textbf{Motion Quality and Text Alignment.} We follow the MoMask~\citep{guo2024momask} evaluation protocol. Fr\'echet Inception Distance (FID) compares generated and reference motion feature distributions. R-Precision at rank 1 (R@1) and multimodal distance (MM-Dist) measure motion--text alignment. Lower FID and MM-Dist and higher R@1 are better. We also report jerk and foot sliding rate (FSR) to measure temporal irregularity and sliding artifacts, following WHAM~\citep{shin2024wham} and OmniControl~\citep{xie2024omnicontrol}; lower values are better for both.
% Before dataset-and-splits revision: \noindent\textbf{Motion Quality and Text Alignment.} We follow the MoMask~\citep{guo2024momask} evaluation protocol. Fr\'echet Inception Distance (FID) compares generated and reference motion feature distributions. R-Precision at rank 1 (R@1) and multimodal distance (MM-Dist) measure motion--text alignment. Lower FID and MM-Dist and higher R@1 are better. We report jerk and foot sliding rate (FSR), following WHAM~\citep{shin2024wham} and OmniControl~\citep{xie2024omnicontrol}. In Table~\ref{tab:comparison}, jerk uses each output's native frame rate and is reported in $10\,\mathrm{m/s^3}$. FSR is the percentage of adjacent frame pairs with sliding on either contacting foot, after resampling to 20 fps. Lower values are better. The ablations use a separate protocol, specified below.
\noindent\textbf{Motion Quality and Text Alignment.} Following MoMask~\citep{guo2024momask}, we report Fr\'echet Inception Distance (FID) for motion distributions, and R-Precision at rank 1 (R@1) and multimodal distance (MM-Dist) for motion--text alignment. Jerk and foot sliding rate (FSR) follow WHAM~\citep{shin2024wham} and OmniControl~\citep{xie2024omnicontrol}. In Table~\ref{tab:comparison}, jerk uses native frame rates and units of $10\,\mathrm{m/s^3}$; FSR measures the percentage of adjacent frame pairs with sliding on either contacting foot after resampling to 20 fps. Ablation protocols differ (Section~\ref{sec:ablation}); table arrows indicate the preferred direction.

% Before citation completion: \noindent\textbf{Scene Interaction.} Following OMOMO, we report Contact-F1. We render generated motions in the corresponding 3D scenes and use Qwen3-VL to assess whether the requested interaction is completed. Interaction Success Rate (ISR) is the fraction of test cases judged successful. These measures complement motion quality: a smooth motion can match the action description while failing to reach the intended object.
\noindent\textbf{Scene Interaction.} Following OMOMO~\citep{li2023omomo}, we report Contact-F1. We render generated motions in the corresponding 3D scenes and use Qwen3-VL~\citep{bai2025qwen3} to assess whether the requested interaction is completed. Interaction Success Rate (ISR) is the fraction of test cases judged successful. These measures complement motion quality: a smooth motion can match the action description while failing to reach the intended object.

% Before citation completion: \noindent\textbf{Video--Motion Consistency.} For joint generation methods, we recover motion from the generated video using CameraHMR and report PoseError against the directly generated motion. This measures agreement between the two outputs through a motion-recovery model, rather than error against a ground-truth 3D trajectory. We therefore interpret PoseError together with text alignment and scene interaction scores.
% Before metric-protocol clarification: \noindent\textbf{Video--Motion Consistency.} For joint generation methods, we recover motion from the generated video using CameraHMR~\citep{patel2025camerahmr} and report PoseError against the directly generated motion. This measures agreement between the two outputs through a motion-recovery model, rather than error against a ground-truth 3D trajectory. We therefore interpret PoseError together with text alignment and scene interaction scores.
\noindent\textbf{Video--Motion Consistency.} For joint generation methods, we recover motion from the generated video using CameraHMR~\citep{patel2025camerahmr} and compare it with the directly generated motion. PoseError is the pelvis-relative mean Euclidean distance over 22 corresponding body joints, in millimeters. We exclude invalid CameraHMR frames and average within each sequence, then across sequences. This measures cross-modal pose agreement, rather than ground-truth 3D accuracy or agreement of the root trajectories.

% Before dataset-and-splits revision: \noindent\textbf{Inference Efficiency.} We report mean inference time in seconds. World2Motion results in the main comparison use the post-processing described in Section~\ref{sec:method}. All ablation results are obtained without the Savitzky--Golay filter so that the effects of the motion representation and noise schedule can be assessed directly.
% Before experiment citation revision: \noindent\textbf{Inference Efficiency.} We report mean inference time in seconds. Main-comparison World2Motion outputs use Savitzky--Golay post-processing (Section~\ref{sec:method}); all ablations disable this filter.
\noindent\textbf{Inference Efficiency.} We report mean inference time in seconds. Main-comparison World2Motion outputs use Savitzky--Golay~\citep{savitzky1964smoothing} post-processing (Section~\ref{sec:method}); all ablations disable this filter.

% Previous quantitative preamble preserved; metrics are defined in Evaluation and the table caption.
% Table~\ref{tab:comparison} compares motion quality, scene interaction, video--motion consistency, and inference time across generation pipelines. We examine these aspects separately because the methods differ in their input requirements and whether they produce both video and motion.
% Continue after the setup without creating a mostly empty page.
\subsection{Quantitative Comparisons}
\begin{table}[!htbp]
\centering
\small
% Before metric-protocol clarification: \caption{Quantitative comparison on the multi-source test set. FID measures motion distribution quality; R@1 and MM-Dist measure text alignment; Contact-F1 and ISR evaluate scene interaction. Jerk and FSR measure motion stability, and PoseError measures video--motion agreement. Time is mean inference time in seconds. Bold denotes the best reported value per column; dashes indicate unreported scores.}
\caption{Quantitative comparison on 100 interaction cases. FID measures motion distribution quality; R@1 and MM-Dist measure text alignment; Contact-F1 and ISR evaluate scene interaction. Jerk and FSR measure motion stability, and PoseError measures video--motion agreement. Time is mean inference time in seconds. Bold denotes the best reported value per column; dashes indicate unreported scores.}
\label{tab:comparison}
\setlength{\tabcolsep}{2pt}
\fitwidth{\textwidth}{%
\begin{tabular}{l c c c c c c c c c}
\toprule
% Before metric-protocol clarification: Method & FID$\downarrow$ & R@1$\uparrow$ & MM-Dist$\downarrow$ & Contact-F1$\uparrow$ & ISR$\uparrow$ & Jerk$\downarrow$ & FSR$\downarrow$ & PoseError$\downarrow$ & Time (s)$\downarrow$\\
Method & FID$\downarrow$ & R@1$\uparrow$ & MM-Dist$\downarrow$ & Contact-F1$\uparrow$ & ISR$\uparrow$ & \begin{tabular}[c]{@{}c@{}}Jerk$\downarrow$\\($10\,\mathrm{m/s^3}$)\end{tabular} & \begin{tabular}[c]{@{}c@{}}FSR$\downarrow$\\(\%)\end{tabular} & \begin{tabular}[c]{@{}c@{}}PoseError$\downarrow$\\(mm)\end{tabular} & Time (s)$\downarrow$\\
\midrule
\multicolumn{10}{l}{\textbf{Text-to-motion}}\\
% Before citation completion: HY-Motion & 0.46 & 0.10 & 1.21 & 0.09 & 0.51 & 3.43 & \textbf{1.25} & -- & 2.83\\
HY-Motion~\citep{hymotion2025} & 0.46 & 0.10 & 1.21 & 0.09 & 0.51 & 3.43 & \textbf{1.25} & -- & 2.83\\
% Before citation completion: Kimodo & 0.66 & 0.05 & 1.30 & 0.03 & 0.23 & 35.14 & 6.07 & -- & \textbf{1.93}\\
Kimodo~\citep{Kimodo2026} & 0.66 & 0.05 & 1.30 & 0.03 & 0.23 & 35.14 & 6.07 & -- & \textbf{1.93}\\
% Before citation completion: Go-to-Zero & 0.51 & 0.04 & 1.35 & 0.07 & 0.36 & 17.50 & 3.42 & -- & 2.57\\
Go-to-Zero~\citep{fan2025gotozero} & 0.51 & 0.04 & 1.35 & 0.07 & 0.36 & 17.50 & 3.42 & -- & 2.57\\
\midrule
\multicolumn{10}{l}{\textbf{Human--object interaction}}\\
% Before citation completion: HOIFHLI & 0.82 & 0.02 & 1.37 & 0.09 & 0.12 & 10.62 & 6.28 & -- & 13.94\\
HOIFHLI~\citep{wu2025hoifhli} & 0.82 & 0.02 & 1.37 & 0.09 & 0.12 & 10.62 & 6.28 & -- & 13.94\\
% Before citation completion: LIGHT & 0.62 & 0.03 & 1.35 & 0.17 & 0.46 & \textbf{1.37} & 11.60 & -- & 25.88\\
LIGHT~\citep{wang2026light} & 0.62 & 0.03 & 1.35 & 0.17 & 0.46 & \textbf{1.37} & 11.60 & -- & 25.88\\
\midrule
\multicolumn{10}{l}{\textbf{Two-stage video generation and motion recovery}}\\
% Before citation completion: \begin{tabular}[c]{@{}l@{}}MiniMax-H3 +\\CameraHMR\end{tabular} & 0.69 & 0.10 & 1.32 & \textbf{0.24} & \textbf{0.91} & 37.22 & 6.41 & -- & 523.67\\
\begin{tabular}[c]{@{}l@{}}MiniMax-H3~\citep{minimaxh3} +\\CameraHMR~\citep{patel2025camerahmr}\end{tabular} & 0.69 & 0.10 & 1.32 & \textbf{0.24} & \textbf{0.91} & 37.22 & 6.41 & -- & 523.67\\
\midrule
\multicolumn{10}{l}{\textbf{Joint video--motion generation}}\\
% Before citation completion: CoMoVi & 0.71 & 0.12 & 1.25 & 0.09 & 0.69 & 16.10 & 24.45 & 81.28 & 270.45\\
CoMoVi~\citep{zhao2026comovi} & 0.71 & 0.12 & 1.25 & 0.09 & 0.69 & 16.10 & 24.45 & 81.28 & 270.45\\
World2Motion & \textbf{0.27} & \textbf{0.24} & \textbf{1.13} & 0.22 & \textbf{0.91} & 8.35 & 2.80 & \textbf{67.50} & 156.59\\
\bottomrule
\end{tabular}
}
\end{table}

% Before concise result discussion: \noindent\textbf{Motion Quality and Text Alignment.} World2Motion achieves the lowest FID and MM-Dist and the highest R@1 among the compared methods. Relative to HY-Motion, FID decreases from 0.46 to 0.27 and R@1 increases from 0.10 to 0.24. The improvement is also observed against CoMoVi, which produces both video and motion but has an FID of 0.71 and R@1 of 0.12. These results show improved alignment between the generated 3D motion and the action prompt, alongside the generation of a corresponding video.
% Before experiment citation revision: \noindent\textbf{Motion Quality and Text Alignment.} World2Motion achieves the lowest FID and MM-Dist and the highest R@1 among the compared methods. Relative to HY-Motion, FID decreases from 0.46 to 0.27 and R@1 increases from 0.10 to 0.24. It also improves over CoMoVi, whose FID and R@1 are 0.71 and 0.12, respectively.
\noindent\textbf{Motion Quality and Text Alignment.} In Table~\ref{tab:comparison}, World2Motion achieves the lowest FID and MM-Dist and the highest R@1 among the compared methods. Relative to HY-Motion, FID decreases from 0.46 to 0.27 and R@1 increases from 0.10 to 0.24. It also improves over CoMoVi, whose FID and R@1 are 0.71 and 0.12, respectively.

\noindent\textbf{Scene Interaction.} World2Motion reaches an ISR of 0.91, compared with 0.23--0.51 for the evaluated text-to-motion methods and 0.12--0.46 for the human--object interaction methods. Its Contact-F1 of 0.22 also exceeds those of these baselines. The two-stage pipeline achieves the same ISR and a higher Contact-F1 of 0.24. Successful completion and contact accuracy are therefore distinct: matching the overall interaction success of the two-stage method does not imply identical contact quality. The remaining contact gap motivates more precise supervision for human--object interactions.

\noindent\textbf{Motion Stability.} World2Motion has lower jerk and FSR than Kimodo, Go-to-Zero, HOIFHLI, and the two video-based baselines. HY-Motion achieves lower values on both stability metrics, while LIGHT achieves the lowest jerk but has a higher FSR of 11.60. Thus, the best result depends on which aspect of stability is measured. World2Motion improves stability over the compared video-based pipelines, but it does not surpass every motion-only method. We retain both jerk and FSR to distinguish temporal irregularity from sliding artifacts.

\noindent\textbf{Video--Motion Consistency.} Compared with CoMoVi, World2Motion reduces PoseError from 81.28 to 67.50, a 17.0\% reduction. This indicates closer agreement between its generated video and directly predicted 3D motion under the motion-recovery evaluation. It complements the improvements in text alignment and scene interaction, which assess the generated motion against the prompt and environment rather than against the other output modality.

\noindent\textbf{Inference Efficiency.} World2Motion takes 156.59 seconds on average, compared with 523.67 seconds for MiniMax-H3 + CameraHMR and 270.45 seconds for CoMoVi: speedups of approximately $3.34\times$ and $1.73\times$, respectively. The text-to-motion and human--object interaction methods remain faster. Our efficiency gain concerns video-based motion generation. World2Motion provides both a video and an explicit 3D motion sequence, whereas the faster motion-only baselines do not produce the corresponding video. The results support reducing the cost of obtaining these paired outputs while preserving competitive interaction performance.

% Continuous preprint layout: no forced page break.
\subsection{Qualitative Comparisons}

\begin{figure}[!htbp]
\centering
\includegraphics[width=0.96\linewidth]{Fig/pdf/Comparison_v2.pdf}
\caption{Qualitative comparison of scene interactions: (a) lifting a plastic box, (b) walking to a chair and sitting, and (c) walking to a sofa and sitting. Each method shows successive motion states; video-generating methods also show corresponding video frames. Human--object interaction baselines include their predicted object trajectories.}
\label{fig:Comparison}
\end{figure}

Figure~\ref{fig:Comparison} compares actions that depend on both body motion and placement relative to an object. In the sitting examples, producing a plausible sitting pose is insufficient: the character must first approach the intended seat and then align the body with it. The box example also involves coordinating the body with object motion. Viewing the approach and final interaction together helps distinguish action matching from scene-consistent execution. For video-based methods, the paired rows additionally allow the reader to inspect whether the generated video and recovered or directly generated 3D motion describe the same interaction.

% Continuous preprint layout: no forced page break.
\subsection{Ablation Studies}
\label{sec:ablation}

% Before metric-protocol clarification: We evaluate root translation and noise scheduling using jerk and FSR. All ablations disable post-processing. Comparisons within each study keep other settings fixed. The first study compares motion representations; the second changes training and inference shifts independently. This separates the choice of noise schedule from the effect of using it consistently.
% Before condensing discussion to accommodate Limitations: We evaluate root translation and noise scheduling using 10,000-iteration checkpoints on a separate 35-case set (15 demonstration clips and 20 internet-video clips), without post-processing. Table~\ref{tab:comparison} instead uses the 20,000-iteration model with smoothing on 100 interaction cases. Ablation jerk uses 25 fps and is reported in $10^3\,\mathrm{m/s^3}$. Ablation FSR is the fraction of contacting foot-joint/frame pairs that slide, pooled over four foot joints on a normalized skeleton, in the range 0--1. Its denominator differs from the all-frame-pair denominator in Table~\ref{tab:comparison}; the scores are therefore compared only within the ablation protocol. The relative-translation row in Table~\ref{tab:ablation_norm} and the matched $s_m=1$ entry in Table~\ref{tab:ablation_schedule} share the same generated outputs.
All ablations use 10,000-iteration checkpoints and 35 clips (15 demonstration clips and 20 internet-video clips), without post-processing; Table~\ref{tab:comparison} uses the smoothed 20,000-iteration model on 100 interaction cases. Ablation jerk uses 25 fps. FSR pools sliding observations over contacting foot-joint/frame pairs from four foot joints on a normalized skeleton. Unlike Table~\ref{tab:comparison}, its denominator excludes non-contact observations, so the scores are not directly comparable across protocols. The relative-translation row in Table~\ref{tab:ablation_norm} and the matched $s_m=1$ entry in Table~\ref{tab:ablation_schedule} share the same outputs.

% Previous root-translation discussion and table; full table moved to Appendix.
% \noindent\textbf{Root translation.} Representing each root position relative to the initial position reduces jerk from 4.355 to 1.251 (71.3\%) and FSR from 0.766 to 0.457 (40.3\%), compared with absolute translation (Table~\ref{tab:ablation_norm}). Both representations retain root position and 6D joint rotations; the change is the reference position used for translation. The comparison therefore isolates this coordinate choice while keeping the rotation representation and other settings fixed. These results support using relative positions for more stable generation without requiring root velocity as an input.
% 
% \begin{table}[!htbp]
% \centering
% \small
% \setlength{\tabcolsep}{5pt}
% \caption{Root translation ablation without post-processing. Relative positions are measured from the initial root position. Bold marks the lower value in each column.}
% \label{tab:ablation_norm}
% \begin{tabular}{lcc}
% \toprule
% Root translation & Jerk $\downarrow$ & FSR $\downarrow$\\
% \midrule
% Absolute & 4.355 & 0.766\\
% Relative to initial position & \textbf{1.251} & \textbf{0.457}\\
% \bottomrule
% \end{tabular}
% \end{table}
% Previous condensed root-translation discussion preserved.
% \noindent\textbf{Root translation.} With other settings fixed and post-processing disabled, expressing root positions relative to the initial position reduces jerk from 4.355 to 1.251 (71.3\%) and FSR from 0.766 to 0.457 (40.3\%), compared with absolute positions (Appendix Table~\ref{tab:ablation_norm}).
\begin{wraptable}{r}{0.45\textwidth}
\centering
\small
% Before metric-protocol clarification: \caption{Root translation ablation without post-processing. Relative positions use the initial root position as the reference.}
% Before qualified ablation metric labels: \caption{Root translation ablation on 35 cases, without post-processing. Relative positions use the initial root position as the reference.}
\setlength{\belowcaptionskip}{6pt}
\caption{Root translation ablation on 35 cases without post-processing. Raw Jerk uses 25 fps ($10^3\,\mathrm{m/s^3}$). Contact FSR is the sliding fraction among contacting foot-joint/frame pairs on a normalized skeleton (0--1); its definition differs from Table~\ref{tab:comparison}.}
\label{tab:ablation_norm}
\setlength{\tabcolsep}{4pt}
\begin{tabular}{l c c}
\toprule
% Before metric-protocol clarification: Translation & Jerk $\downarrow$ & FSR $\downarrow$\\
% Before qualified ablation metric labels: Translation & \begin{tabular}[c]{@{}c@{}}Jerk $\downarrow$\\($10^3\,\mathrm{m/s^3}$)\end{tabular} & \begin{tabular}[c]{@{}c@{}}FSR $\downarrow$\\(0--1)\end{tabular}\\
Translation & \begin{tabular}[c]{@{}c@{}}Raw Jerk $\downarrow$\\($10^3\,\mathrm{m/s^3}$)\end{tabular} & \begin{tabular}[c]{@{}c@{}}Contact FSR $\downarrow$\\(0--1)\end{tabular}\\
\midrule
Absolute & 4.355 & 0.766\\
Relative & \textbf{1.251} & \textbf{0.457}\\
\bottomrule
\end{tabular}
\end{wraptable}
% Before condensing discussion to accommodate Limitations: \noindent\textbf{Root translation.} Table~\ref{tab:ablation_norm} compares absolute root positions with positions relative to the initial frame. With other settings fixed and post-processing disabled, relative translation reduces jerk from 4.355 to 1.251 (71.3\%) and FSR from 0.766 to 0.457 (40.3\%). Both configurations use the same 6D joint rotations; only the translation reference changes. The improvements therefore support this coordinate choice for temporal stability, without introducing root velocity features or relying on smoothing.
\noindent\textbf{Root translation.} With other settings fixed, relative root translation reduces jerk from 4.355 to 1.251 (71.3\%) and FSR from 0.766 to 0.457 (40.3\%) compared with absolute translation (Table~\ref{tab:ablation_norm}). Both configurations use the same 6D joint-rotation representation. These results support the relative coordinate choice for temporal stability in the tested setting, without root velocity features or post-processing.

% Before condensing discussion to accommodate Limitations: \noindent\textbf{Shift-decoupled noise schedule.} Table~\ref{tab:ablation_schedule} separates two questions: which shifts to use, and whether the training and inference shifts should match. With matched shifts, $s_m=1$ achieves lower jerk and FSR than $s_m=10$ or $15$. Compared with matched $s_m=10$, matched $s_m=1$ reduces jerk by 23.5\% and FSR by 28.0\%. For each training configuration, matching the inference shift also gives the lowest values on both metrics. For example, after training with $s_m=1$, changing only the inference shift to $10$ increases jerk from 1.251 to 8.003. These observations support modality-specific shifts and consistent training and inference noise relationships. They do not require an additional smoothing step, since the filter is disabled throughout this comparison.
\noindent\textbf{Shift-decoupled noise schedule.} In Table~\ref{tab:ablation_schedule}, matched training and inference shifts of $s_m=1$ reduce jerk and FSR by 23.5\% and 28.0\% relative to matched $s_m=10$. Matching shifts gives the lowest values within every row. After training with $s_m=1$, changing only the inference shift to $10$ increases jerk from 1.251 to 8.003. These results support modality-specific shifts and training--inference consistency for the tested stability metrics.

% Previous nine-row schedule table preserved.
% \begin{table}[!htbp]
% \centering
% \small
% \setlength{\tabcolsep}{5pt}
% \caption{Noise-schedule ablation without post-processing. Video shift $s_v=10$ is fixed in training and inference. Groups vary the training motion shift; bold marks the best result within each group.}
% \label{tab:ablation_schedule}
% \begin{tabular}{cccc}
% \toprule
% Training $s_m$ & Inference $s_m$ & Jerk $\downarrow$ & FSR $\downarrow$\\
% \midrule
% \multirow{3}{*}{1} & 1 & \textbf{1.251} & \textbf{0.457}\\
%  & 10 & 8.003 & 0.986\\
%  & 15 & 9.703 & 1.000\\
% \midrule
% \multirow{3}{*}{10} & 1 & 7.021 & 0.985\\
%  & 10 & \textbf{1.635} & \textbf{0.635}\\
%  & 15 & 2.557 & 0.830\\
% \midrule
% \multirow{3}{*}{15} & 1 & 8.413 & 0.981\\
%  & 10 & 2.127 & 0.700\\
%  & 15 & \textbf{1.809} & \textbf{0.649}\\
% \bottomrule
% \end{tabular}
% \end{table}
% Previous two-block schedule layout preserved.
% \begin{table}[!htbp]
% \centering
% \small
% \caption{Noise-schedule ablation without post-processing. Video shift $s_v=10$ is fixed in training and inference. Bold marks the lowest value in each training row.}
% \label{tab:ablation_schedule}
% \setlength{\tabcolsep}{5pt}
% \begin{minipage}[t]{0.48\linewidth}
% \centering
% \textbf{Jerk $\downarrow$}\par\smallskip
% \begin{tabular}{c c c c}
% \toprule
%  & \multicolumn{3}{c}{Inference $s_m$}\\
% \cmidrule(lr){2-4}
% Training $s_m$ & 1 & 10 & 15\\
% \midrule
% 1 & \textbf{1.251} & 8.003 & 9.703\\
% 10 & 7.021 & \textbf{1.635} & 2.557\\
% 15 & 8.413 & 2.127 & \textbf{1.809}\\
% \bottomrule
% \end{tabular}
% \end{minipage}\hfill
% \begin{minipage}[t]{0.48\linewidth}
% \centering
% \textbf{FSR $\downarrow$}\par\smallskip
% \begin{tabular}{c c c c}
% \toprule
%  & \multicolumn{3}{c}{Inference $s_m$}\\
% \cmidrule(lr){2-4}
% Training $s_m$ & 1 & 10 & 15\\
% \midrule
% 1 & \textbf{0.457} & 0.986 & 1.000\\
% 10 & 0.985 & \textbf{0.635} & 0.830\\
% 15 & 0.981 & 0.700 & \textbf{0.649}\\
% \bottomrule
% \end{tabular}
% \end{minipage}
% \end{table}
\begin{table}[!htbp]
\centering
\small
% Before metric-protocol clarification: \caption{Noise-schedule ablation without post-processing. Columns within each metric group specify the inference motion shift $s_m$. Video shift $s_v=10$ is fixed throughout. Bold marks the lowest value for each metric within each row.}
% Before qualified ablation metric labels: \caption{Noise-schedule ablation on the same 35 cases as Table~\ref{tab:ablation_norm}, without post-processing. Columns within each metric group specify the inference motion shift $s_m$. Video shift $s_v=10$ is fixed throughout. Bold marks the lowest value for each metric within each row.}
\setlength{\belowcaptionskip}{6pt}
\caption{Noise-schedule ablation on the same 35 cases, using Raw Jerk and contact-conditioned FSR as defined in Table~\ref{tab:ablation_norm}. These unfiltered ablation scores are not directly comparable to Table~\ref{tab:comparison}. Columns give the inference motion shift $s_m$; $s_v=10$ is fixed. Bold marks the lowest value per metric within each row.}
\label{tab:ablation_schedule}
\setlength{\tabcolsep}{7pt}
\begin{tabular}{c c c c c c c}
\toprule
% Before metric-protocol clarification:  & \multicolumn{3}{c}{Jerk $\downarrow$} & \multicolumn{3}{c}{FSR $\downarrow$}\\
% Before qualified ablation metric labels:  & \multicolumn{3}{c}{Jerk ($10^3\,\mathrm{m/s^3}$) $\downarrow$} & \multicolumn{3}{c}{FSR (0--1) $\downarrow$}\\
 & \multicolumn{3}{c}{Raw Jerk ($10^3\,\mathrm{m/s^3}$) $\downarrow$} & \multicolumn{3}{c}{Contact FSR (0--1) $\downarrow$}\\
\cmidrule(lr){2-4}\cmidrule(lr){5-7}
Training $s_m$ & 1 & 10 & 15 & 1 & 10 & 15\\
\midrule
1 & \textbf{1.251} & 8.003 & 9.703 & \textbf{0.457} & 0.986 & 1.000\\
10 & 7.021 & \textbf{1.635} & 2.557 & 0.985 & \textbf{0.635} & 0.830\\
15 & 8.413 & 2.127 & \textbf{1.809} & 0.981 & 0.700 & \textbf{0.649}\\
\bottomrule
\end{tabular}
\end{table}
% Let floats share the following text page.